%%This is a template for PH1140 documents for RHUL physics. It is based on IOP journal formatting, but has been brought up-to-date with biblatex by ACK (2017). 


\documentclass[12pt]{iopart}
\usepackage{graphicx} %includes figure handling
\usepackage[sorting=none]{biblatex} %includes bibliography tools

\defbibheading{iophead}[\refname]{\section*{#1}}% don't touch this!

\bibliography{references.bib} %this is the name of the file where the biblatex entries are written down

\usepackage[a4paper, left=3.5cm,right=3.5cm]{geometry}

\usepackage{tikz}
\usetikzlibrary{calc}
\usepackage{graphicx}
\usepackage{newtxtext}
\usepackage{float}
\usepackage{comment}

\begin{document}

% \title[Your Name Here]{Title of the Report}

% \author{Your Name Here \\ \today} %\today can be changed to a fixed date if desired


% To add this template to the main.tex file, just add the command "% To add this template to the main.tex file, just add the command "% To add this template to the main.tex file, just add the command "\include{titlePageSLU} after "\begin{docuemnt}" in the main.tex file


% In this segment, enter the desired data to be shown at the title page
\newcommand{\thesisAuthor}{SH Khoza}
\newcommand{\thesisNumber}{214651459}
\newcommand{\thesisTitle}{Interesting Thesis Title}
\newcommand{\thesisSubTitle}{little bit more descriptive}
\newcommand{\thesisTitleTranslated}{Translated Headline}
\newcommand{\thesisDegree}{Master thesis project}
\newcommand{\university}{Swedish University of Argicultural Science, SLU}
\newcommand{\credits}{30 hp}
\newcommand{\faculty}{Faculty of blablabalba}
\newcommand{\thesisPlaceDate}{Place of puplication, 2020}
\newcommand{\company}{Company name}

%------------------------------------------------------------------------------
\begin{titlepage}
\thispagestyle{empty}
\myfont

% Use this line of code if both SLU loggo and company/other institution loggo is desired. The positions are possible to change with the \hspace and \vspace syntax.
\begin{figure} [H]
\vspace{-3cm}
 \centering
\begin{minipage}[t]{.45\linewidth}
  \raggedright
  % Upload and include SLU loggo here:
  \hspace*{-2cm}\includegraphics[width=\linewidth]{slu_logo_webb.png}
  
\end{minipage}%
  \begin{minipage}[t]{.45\linewidth}
  \vspace{-3.3cm}
 \raggedleft
% Upload and include other loggo here (loggo of wikipedia is used as an example):
 \hspace*{1cm}\includegraphics[width =0.5\textwidth]{wikilogo.png} \hspace*{-1cm}
 
\end{minipage}
\end{figure}


% If only the logo for SLU is desired, delete "\begin{comment}" and "\end{comment}" to use this line of code:
\begin{comment}
\begin{figure}
\vspace{-3cm}
    \hspace*{-2cm}\includegraphics[width = 0.3\textwidth]{slu_logo_webb.png}\hspace*{-2cm}
\end{figure}
\end{comment}


% Adds background picture. Delete code if no background picture is wanted.
\begin{tikzpicture}[overlay, remember picture]
\node[anchor=south west, 
      xshift=-0.2cm, 
      yshift=-0.2cm] 
     at (current page.south west)
     {\includegraphics[width = 1.8\textwidth, height = 9cm]{background.png}}; 
\end{tikzpicture}


\vspace{1cm}
\par
\noindent
\Huge
\textbf{\thesisTitle}
\vspace{0.2cm}
\LARGE
\par
\noindent
- \thesisSubTitle\\
\rule[0.3cm]{\linewidth}{2pt}
\Large

% Delete this line if no translation is desired
\noindent
\textit{\thesisTitleTranslated}

\vspace{2cm}
\noindent
\LARGE
\thesisAuthor\\
\thesisNumber\\
\vspace{4 cm}
\small
\par \noindent
\thesisDegree $\cdot$ \credits
\par \noindent
\university
\par \noindent
\faculty
\par \noindent
\company
\par \noindent
\thesisPlaceDate

\end{titlepage}
 after "\begin{docuemnt}" in the main.tex file


% In this segment, enter the desired data to be shown at the title page
\newcommand{\thesisAuthor}{SH Khoza}
\newcommand{\thesisNumber}{214651459}
\newcommand{\thesisTitle}{Interesting Thesis Title}
\newcommand{\thesisSubTitle}{little bit more descriptive}
\newcommand{\thesisTitleTranslated}{Translated Headline}
\newcommand{\thesisDegree}{Master thesis project}
\newcommand{\university}{Swedish University of Argicultural Science, SLU}
\newcommand{\credits}{30 hp}
\newcommand{\faculty}{Faculty of blablabalba}
\newcommand{\thesisPlaceDate}{Place of puplication, 2020}
\newcommand{\company}{Company name}

%------------------------------------------------------------------------------
\begin{titlepage}
\thispagestyle{empty}
\myfont

% Use this line of code if both SLU loggo and company/other institution loggo is desired. The positions are possible to change with the \hspace and \vspace syntax.
\begin{figure} [H]
\vspace{-3cm}
 \centering
\begin{minipage}[t]{.45\linewidth}
  \raggedright
  % Upload and include SLU loggo here:
  \hspace*{-2cm}\includegraphics[width=\linewidth]{slu_logo_webb.png}
  
\end{minipage}%
  \begin{minipage}[t]{.45\linewidth}
  \vspace{-3.3cm}
 \raggedleft
% Upload and include other loggo here (loggo of wikipedia is used as an example):
 \hspace*{1cm}\includegraphics[width =0.5\textwidth]{wikilogo.png} \hspace*{-1cm}
 
\end{minipage}
\end{figure}


% If only the logo for SLU is desired, delete "\begin{comment}" and "\end{comment}" to use this line of code:
\begin{comment}
\begin{figure}
\vspace{-3cm}
    \hspace*{-2cm}\includegraphics[width = 0.3\textwidth]{slu_logo_webb.png}\hspace*{-2cm}
\end{figure}
\end{comment}


% Adds background picture. Delete code if no background picture is wanted.
\begin{tikzpicture}[overlay, remember picture]
\node[anchor=south west, 
      xshift=-0.2cm, 
      yshift=-0.2cm] 
     at (current page.south west)
     {\includegraphics[width = 1.8\textwidth, height = 9cm]{background.png}}; 
\end{tikzpicture}


\vspace{1cm}
\par
\noindent
\Huge
\textbf{\thesisTitle}
\vspace{0.2cm}
\LARGE
\par
\noindent
- \thesisSubTitle\\
\rule[0.3cm]{\linewidth}{2pt}
\Large

% Delete this line if no translation is desired
\noindent
\textit{\thesisTitleTranslated}

\vspace{2cm}
\noindent
\LARGE
\thesisAuthor\\
\thesisNumber\\
\vspace{4 cm}
\small
\par \noindent
\thesisDegree $\cdot$ \credits
\par \noindent
\university
\par \noindent
\faculty
\par \noindent
\company
\par \noindent
\thesisPlaceDate

\end{titlepage}
 after "\begin{docuemnt}" in the main.tex file


% In this segment, enter the desired data to be shown at the title page
\newcommand{\thesisAuthor}{SH Khoza}
\newcommand{\thesisNumber}{214651459}
\newcommand{\thesisTitle}{Interesting Thesis Title}
\newcommand{\thesisSubTitle}{little bit more descriptive}
\newcommand{\thesisTitleTranslated}{Translated Headline}
\newcommand{\thesisDegree}{Master thesis project}
\newcommand{\university}{Swedish University of Argicultural Science, SLU}
\newcommand{\credits}{30 hp}
\newcommand{\faculty}{Faculty of blablabalba}
\newcommand{\thesisPlaceDate}{Place of puplication, 2020}
\newcommand{\company}{Company name}

%------------------------------------------------------------------------------
\begin{titlepage}
\thispagestyle{empty}
\myfont

% Use this line of code if both SLU loggo and company/other institution loggo is desired. The positions are possible to change with the \hspace and \vspace syntax.
\begin{figure} [H]
\vspace{-3cm}
 \centering
\begin{minipage}[t]{.45\linewidth}
  \raggedright
  % Upload and include SLU loggo here:
  \hspace*{-2cm}\includegraphics[width=\linewidth]{slu_logo_webb.png}
  
\end{minipage}%
  \begin{minipage}[t]{.45\linewidth}
  \vspace{-3.3cm}
 \raggedleft
% Upload and include other loggo here (loggo of wikipedia is used as an example):
 \hspace*{1cm}\includegraphics[width =0.5\textwidth]{wikilogo.png} \hspace*{-1cm}
 
\end{minipage}
\end{figure}


% If only the logo for SLU is desired, delete "\begin{comment}" and "\end{comment}" to use this line of code:
\begin{comment}
\begin{figure}
\vspace{-3cm}
    \hspace*{-2cm}\includegraphics[width = 0.3\textwidth]{slu_logo_webb.png}\hspace*{-2cm}
\end{figure}
\end{comment}


% Adds background picture. Delete code if no background picture is wanted.
\begin{tikzpicture}[overlay, remember picture]
\node[anchor=south west, 
      xshift=-0.2cm, 
      yshift=-0.2cm] 
     at (current page.south west)
     {\includegraphics[width = 1.8\textwidth, height = 9cm]{background.png}}; 
\end{tikzpicture}


\vspace{1cm}
\par
\noindent
\Huge
\textbf{\thesisTitle}
\vspace{0.2cm}
\LARGE
\par
\noindent
- \thesisSubTitle\\
\rule[0.3cm]{\linewidth}{2pt}
\Large

% Delete this line if no translation is desired
\noindent
\textit{\thesisTitleTranslated}

\vspace{2cm}
\noindent
\LARGE
\thesisAuthor\\
\thesisNumber\\
\vspace{4 cm}
\small
\par \noindent
\thesisDegree $\cdot$ \credits
\par \noindent
\university
\par \noindent
\faculty
\par \noindent
\company
\par \noindent
\thesisPlaceDate

\end{titlepage}


\begin{abstract}
The abstract goes here, ensuring that the environment is opened and closed with \texttt{abstract}.

\end{abstract}

\section{Sections}

Sections should be delineated by \texttt{\textbackslash section}. Avoid unnecessary subdivision of text.

\subsection{Subsections}

Subsections should be delineated by \texttt{\textbackslash subsection}.

\subsubsection{Subsubsections}

Subsubsections should be delineated by \texttt{\textbackslash subsubsection}.

\section{Labels \label{sec:label}}

Almost any part of a \LaTeX document can be labeled by using the \texttt{\textbackslash label} command, and then refered to. For example, this is Section~\ref{sec:label}.

\section{Equations}

Equations can be typeset with the \texttt{\textbackslash equation} environment.

\begin{equation}
V=\frac{nkT}{P}
\label{eq:gaslaw}
\end{equation}

If you need to talk about variables within equations, you can use \$...\$. For example, on the left hand side of Eq.~\ref{eq:gaslaw}, $V$ stands for volume. Note the use of a label to refer to the equation!

\section{Figures}

Figures are important. They can be included as .pdf, .eps, .jpg, or .png.
However, .pdf and .eps are recommended for plots, as they are vector graphics
formats, which have better scaling properties to various sizes. 

\begin{figure}
  \centering
  \includegraphics[width=0.6\textwidth]{logo-large-london-cmyk.pdf}\\
  \caption{Figure captions go underneath the figure. This is the Royal Holloway logo.}
  \label{fig:logo}
\end{figure}


\section{Tables}

Tables can be complicated beasts! \LaTeX will prevent tables from spanning two pages, but you should help it as much as you can by intelligent formatting. Table~\ref{tab:somenumber} shows some totally bogus data!


\begin{table}
  \centering
  \caption{Captions of tables go on top of tables. This table contains some made up numbers.}
  \label{tab:somenumber}
\begin{tabular}{|c|c|c|}
\hline
Length  & Voltage  & Angle \\
(mm) & (V) & (degrees)\\
\hline 
5$\pm$1 & 10$\pm$ 0.1 & 0 $\pm$ 0.01\\
7 & 14 & 3.5\\
10 & 20 & 5\\
\hline
\end{tabular}
\end{table}

\section{Referencing}

Biblatex is the best way to reference in \LaTeX. References go in a separate .bib file, and each has a label associated with it, so you can cite, for example, a journal article~\cite{ahmad2002direct}, a book~\cite{universityphysics}, a website~\cite{Biblatex}, or a technical manual~\cite{PASCO}.


\printbibliography[heading=iophead]

\end{document}

